\documentclass[12pt,a4paper]{article}

\usepackage[T1]{fontenc}
\usepackage{lmodern}
\input{glyphtounicode}
\pdfgentounicode=1

\usepackage{amsmath, amsthm, amssymb}
\usepackage[hidelinks]{hyperref}
\hypersetup{
 pdftitle={Toward Coefficient Isotropy from Character Uniformity: A Stability Lemma and Its Missing Geometric Inputs},
 pdfauthor={J\'er\^ome Beau},
 pdfkeywords={Heisenberg group; Weil representation; spectral universality; Gram-Schmidt rank increment;
 complex conjugation; su(2)-invariant forms; Schur's lemma; coefficient isotropy; emergent geometry; Cosmochrony}
}
\usepackage{geometry}
\usepackage{booktabs}
\geometry{margin=2.5cm}

\newtheorem{theorem}{Theorem}[section]
\newtheorem{lemma}[theorem]{Lemma}
\newtheorem{proposition}[theorem]{Proposition}
\newtheorem{corollary}[theorem]{Corollary}
\theoremstyle{definition}
\newtheorem{definition}[theorem]{Definition}
\newtheorem*{hypo}{Input}
\theoremstyle{remark}
\newtheorem{remark}[theorem]{Remark}

\DeclareMathOperator{\Heis}{Heis}
\DeclareMathOperator{\Sym}{Sym}
\DeclareMathOperator{\Var}{Var}
\DeclareMathOperator{\spn}{span}
\newcommand{\Heff}{H_{\mathrm{eff}}}
\newcommand{\Csu}{C_{\mathfrak{su}(2)}}
\newcommand{\spair}{\sigma_{\mathrm{pair}}}
\newcommand{\dpair}{\delta_{\mathrm{pair}}}
\newcommand{\Bq}{\mathfrak{B}_q}
\newcommand{\Nq}{N_q}
\newcommand{\Cpm}{\mathcal{C}_{\pm}}
\newcommand{\Czero}{\mathcal{C}_{0}}

\bibliographystyle{plain}
\usepackage{doi}

\begin{document}

  \title{%
    Toward Coefficient Isotropy from Character Uniformity:\\
    A Stability Lemma and Its Missing Geometric Inputs\\[0.4em]
    {\large Q10 of the Cosmochrony Spectral Geometry Programme}%
  }
  \author{J\'er\^ome Beau\\
  \small Independent Researcher\\
  \small \href{mailto:jerome.beau@cosmochrony.org}{jerome.beau@cosmochrony.org}
  }
  \date{2026}
  \maketitle

  \begin{abstract}
    An isotropic spatial geometry reconstructed from the admissible fibre of the Cosmochrony programme would need the
    horizontal and central coefficients of a positive three-dimensional spatial form to agree, and would need their
    common value.
    This paper isolates what uniformity of the O-series fingerprint observable across characters can contribute.
    The observable is a Gram--Schmidt rank increment per breadth-first shell; the hypothesis [U] asks that, on a
    specified block family and depth window, it stay within a relative distance $\varepsilon$ of a common positive
    profile.
    We prove a sharp lemma: two averages of such data with a common depth weighting have a ratio within
    $2\varepsilon/(1-\varepsilon)$ of one.
    This becomes a statement about geometric coefficients only through four inputs named and not supplied here: a
    sector assignment of blocks to spin weights, a linear response law, a positive rank-three target form, and an
    independent absolute scale.
    Under them the coefficient ratio obeys the same bound, and the absolute coefficient follows only from the
    supplied scale.
    The $\mathfrak{su}(2)$ Casimir acts as 2 on the spin-one module, but invariance fixes an invariant form only up
    to a positive factor, so the value 2 determines no coefficient.
    Exact conjugation parity holds for conjugately matched fingerprint data with an identical normalisation, which
    the published pair campaign, sampling its blocks independently, does not meet.
    Parity and concentrated pair exponents do not imply [U], and a three-dimensional neutral sector and rank
    stability constrain none of the amplitudes that [U] bounds.
    The available compression entries are Rayleigh-quotient identities of the discrete Weil Laplacian, not
    coefficients; no coefficient value is derived.
    Interpretive outlook: along this route, relative isotropy would be inherited from the uniformity of the
    admissible fibre across characters, while the absolute scale of emergent space would remain an independent
    datum rather than a consequence of representation theory.
  \end{abstract}

  \bigskip
  \noindent\textbf{Keywords.}
  Heisenberg group; Weil representation; spectral universality; Gram--Schmidt rank increment; complex conjugation;
  su(2)-invariant forms; Schur's lemma; coefficient isotropy; emergent geometry; Cosmochrony.

  \noindent\textbf{MSC 2020.}
  22E25, 81R05, 53B30, 15A63.

  \clearpage
  \tableofcontents
  \clearpage

  \section{Introduction}

    \subsection{The isotropy question and what is available}

      Paper Q5b~\cite{Beau2026q5b} extracts, conditionally on its spatial limit hypothesis [H-L] and on the lifting
      hypothesis [H-lift], a leading-order effective co-metric $\mathrm{diag}(-A_\tau, A_H, A_H, 0)$ on
      $\mathbb{R}_\tau \times \Heis_3(\mathbb{R})$ (Q5b Theorem~5.2).
      Its spatial principal form has rank two everywhere, and the central slot is empty.
      Paper Q9~\cite{Beau2026q9} proves a kinetic Mosco statement under explicit inputs and states that its results
      do not discharge [H-lift], which Q5b records as open.
      Paper Q7~\cite{Beau2026q7} states what an equivariant identification between the three-dimensional admissible
      space $\Heff$ and a three-dimensional spatial sector would require.
      A positive rank-three spatial form must come from a new construction beyond Q5b's operator (Q7
      Conjecture~4.7); if one is supplied together with the compatibility conditions of Q7, its pullback to the
      spin-one module has a graded form with two positive coefficients, a horizontal $A_H$ and a central $A_Z$
      (Q7 Lemma~4.5).
      No value of $A_H$ or $A_Z$ is established: Q5b records the coefficient values as open (its problem Q5b-O3),
      Theorem~\ref{thm:main} below derives no value of $A_H$, and Q8~\cite{BeauQ8} shows that the Heisenberg
      commutator supplies neither a constant positive central
      coefficient nor a rank-three symbol, and that $\mathfrak{su}(2)$-invariance leaves the common value of an
      isotropic form free.

      The present paper asks a narrower question: what can uniformity of the O-series fingerprint observable across
      characters contribute to the equality $A_H = A_Z$, and what else would be needed to reach a value?

    \subsection{Main result}

      The inputs [W], [R], [T] and [S] are stated in Section~\ref{sec:setup}; [U] is
      Definition~\ref{def:universality}.

      \begin{theorem}[Conditional coefficient isotropy]\label{thm:main}
        Let $q$ be an odd prime.
        Assume \textup{[U]} with rate $\varepsilon(q) < 1$ on a block family $\Bq$ and a depth window $\Nq$, the
        sector assignment \textup{[W]}, the response law \textup{[R]} and the supplied target \textup{[T]}.
        Then
        \begin{equation}\label{eq:main}
          \left|\frac{A_H(q)}{A_Z(q)} - 1\right| \;\leq\; \frac{2\,\varepsilon(q)}{1 - \varepsilon(q)} ,
        \end{equation}
        and $A_H(q)/A_Z(q) \to 1$ along any sequence of primes with $\varepsilon(q) \to 0$.
        If in addition the scale input \textup{[S]} supplies $A_Z(q) = a_q$, independently of $A_H$, then
        $|A_H(q) - a_q| \leq 2\,a_q\,\varepsilon(q)/(1-\varepsilon(q))$, and $A_H(q) \to a_*$ whenever
        $a_q \to a_* > 0$ and $\varepsilon(q) \to 0$.
        No value of $A_H$, $A_Z$ or $a_*$ is derived here.
      \end{theorem}

      The bound~\eqref{eq:main} is the sharp bound for the underlying averages (Lemma~\ref{lem:su2-isotropy}).
      The content of Theorem~\ref{thm:main} beyond that lemma lies entirely in the inputs: [R] is the only link
      between fingerprint data and a geometric coefficient, [T] is the only source of a positive central
      coefficient, and [S] is the only source of an absolute value.
      None of [U], [W], [R], [T] or [S] is established in the corpus.
      The $\mathfrak{su}(2)$ Casimir does not replace [S] (Proposition~\ref{prop:casimir-unique}).

    \subsection{Structure of the argument}

      Section~\ref{sec:setup} fixes the observable and states the inputs.
      Section~\ref{sec:universality} states [U], shows that the inputs cited in its support do not imply it, and
      gives the exact conditions for conjugation parity (Lemma~\ref{lem:conjugation}).
      Section~\ref{sec:su2} proves the stability lemma and its conditional coefficient form.
      Section~\ref{sec:casimir} records what $\mathfrak{su}(2)$-invariance does and does not fix, and proves
      Theorem~\ref{thm:main}.
      Section~\ref{sec:numerical} types the available numbers.
      Section~\ref{sec:consequences} states the consequences and the open inputs, and Section~\ref{sec:conclusion}
      concludes.

  \section{Setup and Inputs}\label{sec:setup}

    \subsection{What Q5b and Q7 supply}

      Let $L_{\mathrm{eff}} = -A_\tau\partial_\tau^2 + A_H\,\Delta_H$ be the effective operator of Q5b
      Theorem~5.2~\cite{Beau2026q5b}, conditional on [H-L] and [H-lift], with $\Delta_H = X^2 + Y^2$ the
      sub-Laplacian of $\Heis_3(\mathbb{R})$.
      Its spatial principal symbol has rank two at every point, so it is not congruent to any positive rank-three
      form (Q7, discussion after Conjecture~4.7~\cite{Beau2026q7}); terms of lower differential order do not change
      it.
      Consequently no central coefficient $A_Z$ is defined by $L_{\mathrm{eff}}$.

      Throughout, $V_\rho \simeq \mathbb{C}^2$ is the spin-$\frac12$ module of the supplied carrier of
      O23~\cite{Beau2026a27}, and $\Sym^2(V_\rho)$ is the spin-one module with weight vectors
      $e_+, e_0, e_-$ of $J_3$-weights $+1, 0, -1$.
      The identification of $\Heff$ with $\Sym^2(V_\rho)$ is Q7's Hypothesis [ID]; the adjoint action on the spatial
      target is [ACT]; the rotation and real-structure compatibility is [INV]; the adapted normalisation is [ADAPT]
      (Q7 Section~4 and Remark~4.6).
      None of these is established for a geometric construction.

      \begin{hypo}[{[T]}, supplied target]
        For each $q$ there is a candidate continuum operator $L_{\mathrm{ext}}$ (Q7 Criterion~5.1) whose spatial
        form is positive of rank three, together with [ID], [ACT], [INV] and [ADAPT], such that its pullback along the
        adapted isomorphism $\varphi_0$ to $\Sym^2(V_\rho)$ is $A_H(q)\,Q_\pm + A_Z(q)\,Q_0$ with $A_H(q),
        A_Z(q) > 0$, where $Q_\pm$ and $Q_0$ are
        the projections onto $\spn\{e_+, e_-\}$ and $\spn\{e_0\}$ (Q7 Lemma~4.5(3)).
      \end{hypo}

      By Q7 Lemma~4.5(3), the transport along $\lambda\varphi_0$ multiplies both coefficients by $|\lambda|^2$; the
      ratio $A_H/A_Z$ does not depend on that choice, and along $\varphi_0$ the coefficients are the entries of the
      spatial form in the adapted frame.
      Hypothesis [ID] belongs to [T] here because [W] and [R] relate fingerprint data on the admissible space to the
      spin-one module; Q8~\cite{BeauQ8} states the same target input without [ID], since its statements concern only
      the module and the target.

    \subsection{The fingerprint observable}

      We follow the O12/O13 protocol as used in O25~\cite{Beau2026a29}.
      A block is a triple $b = (c_1, c_2, c_3)$ of nonzero residues mod $q$ with $c_1 + c_2 + c_3 \not\equiv 0$;
      its first entry is its character.
      For each depth $n$ let $S_n$ be the $n$-th breadth-first shell of the Cayley graph of
      $\Heis_3(\mathbb{Z}/q\mathbb{Z})$ for the fixed generating set; the shells do not depend on the block.
      The observable is
      \begin{equation}\label{eq:sigmac}
        \sigma_b(n) := \frac{\delta r_n(b)}{|S_n|},
      \end{equation}
      where $\delta r_n(b)$ is the number of Weil fingerprint vectors of block $b$ in the shell $S_n$ that are
      linearly independent of the Gram--Schmidt span built over the shells $0, \dots, n-1$ (O25, Section~2.1).
      We read this count as the increase of the dimension of the span from shell $n-1$ to shell $n$, which is the
      value of the sequential count in which each vector of $S_n$ is tested against everything processed before it.
      It is a rank increment, an integer, not a spectral quantity; U1~\cite{Beau2026u1} records the same typing.
      Independence is understood in exact arithmetic.
      The pair observable of a conjugate pair is
      \begin{equation}\label{eq:sigmapair}
        \spair(n) := \sigma_b(n)\,\sigma_{b'}(n),
      \end{equation}
      where $b$ has character $c$ and $b'$ has character $q - c$; in O25 the remaining entries of $b$ and $b'$ are
      sampled independently and uniformly (O25, Section~2.2).
      The pair exponent $\dpair$ is minus the least-squares slope of $\log \spair$ against $\log(n+1)$ over an
      auto-calibrated fitting window.

    \subsection{Sector assignment, response law and scale}

      \begin{hypo}[{[W]}, sector assignment]
        For each $q$ there are nonempty disjoint finite sets of blocks $\Cpm$ and $\Czero$, assigned respectively to
        the spin weights $\pm 1$ and to the spin weight $0$ of $\Sym^2(V_\rho)$.
      \end{hypo}

      No source constructs such an assignment.
      Q7 Remark~4.2 relates spin weights to Carnot degrees, subspace to subspace; it assigns no sets of characters
      or blocks.

      \begin{hypo}[{[R]}, linear response law]
        For each $q$ there are probability weights $\alpha_H$ on $\Cpm$ and $\alpha_Z$ on $\Czero$ and one probability
        weight $\nu_q$ on the depth window $\Nq$ such that
        \begin{equation}\label{eq:ratio}
          \frac{A_H(q)}{A_Z(q)}
          = \frac{\sum_{b \in \Cpm} \sum_{n \in \Nq} \alpha_H(b)\,\nu_q(n)\,\sigma_b(n)}
                 {\sum_{b \in \Czero} \sum_{n \in \Nq} \alpha_Z(b)\,\nu_q(n)\,\sigma_b(n)},
        \end{equation}
        where $A_H(q)$ and $A_Z(q)$ are the coefficients of~\textup{[T]}.
      \end{hypo}

      Input [R] is the geometric content of the whole construction.
      It is not a consequence of Q7 Criterion~5.1, which is a necessary form test on a supplied continuum operator
      and contains no averaging measure, no depth weighting and no map from discrete compressions to continuum
      coefficients; Q7 states that no source supplies that transfer.
      The common depth weight $\nu_q$ is essential: with different depth weights in the numerator and the
      denominator, the ratio~\eqref{eq:ratio} can stay away from $1$ under exact uniformity.

      \begin{hypo}[{[S]}, independent scale]
        For each $q$ a value $a_q > 0$ of $A_Z(q)$ is supplied by a construction that does not use $A_H$.
      \end{hypo}

      Without [S], Theorem~\ref{thm:main} is a statement about a ratio only: any pair $A_H = A_Z = a$ with $a > 0$ is
      compatible with exact isotropy.

  \section{The Uniformity Hypothesis}\label{sec:universality}

    \subsection{Statement}

      \begin{definition}[Spectral universality, {\upshape [U]}]\label{def:universality}
        Let $\Bq$ be a finite family of blocks containing $\Cpm \cup \Czero$, let $\Nq$ be a finite nonempty depth
        window, and let $\sigma_* \colon \Nq \to (0, \infty)$ be a reference profile.
        The family satisfies \textup{[U]} with rate $\varepsilon(q) \in [0, 1)$ if
        \begin{equation}\label{eq:universality}
          \bigl|\sigma_b(n) - \sigma_*(n)\bigr| \;\leq\; \varepsilon(q)\,\sigma_*(n)
          \qquad \text{for all } b \in \Bq,\ n \in \Nq.
        \end{equation}
        As a limit statement, \textup{[U]} asks that such data exist along a sequence of primes with
        $\varepsilon(q) \to 0$.
      \end{definition}

      No functional form of $\sigma_*$ is assumed.
      Positivity of $\sigma_*$ on the window is part of the definition: at depths where the Gram--Schmidt span has
      saturated $\mathbb{C}^q$, every $\sigma_b(n)$ vanishes and a relative bound carries no information.
      A formulation indexed by characters alone requires a rule selecting the other two block entries; that rule is
      part of the choice of $\Bq$.

    \subsection{What the cited inputs give}

      \begin{proposition}[Status of the inputs cited for \textup{[U]}]\label{prop:structural-U}
        The following statements hold.
        \begin{enumerate}
          \item[(i)] \textup{(Conjugation.)}
          If a block $\bar b$ is conjugately matched to a block $b$, with an identical normalisation, then
          $\sigma_{\bar b} = \sigma_b$ at every depth (Lemma~\ref{lem:conjugation}).
          O17~\cite{Beau2026a21} proves exact equality of raw redundancy counts in its own scalar Weil model,
          where the count is in fact independent of the block (O17 Theorem~2.1 and Corollary~2.3), and states that
          this model is not shown to coincide with the O12/O13 pipeline.
          O18~\cite{Beau2026a22} recalls that result (O18 Theorem~3.1) and keeps the fibre reading of the pair
          $\{c, q-c\}$ as an open bridge (O18 Remark~3.2).
          O22~\cite{Beau2026a26} uses the conjugation invariance of the pair observable $\spair$, which is symmetric
          in its two factors; invariance of a product under exchange of its factors does not give equality of the
          factors.

          \item[(ii)] \textup{(Supplied carrier.)}
          O23~\cite{Beau2026a27} proves that the traceless anti-Hermitian sector of a supplied two-dimensional
          $SU(2)$ carrier has three independent directions (O23 Theorem~3.1); the carrier selection and the threshold
          $\Sigma_c(n_3) = 3$ are supplied.
          This fixes no amplitude profile of $\sigma_b$.

          \item[(iii)] \textup{(Rank stability.)}
          O24~\cite{Beau2026a28} excludes observable-rank changes induced by vertical fibre structure, conditionally
          on the supplied carrier and fibre hypotheses.
          This bounds no relative deviation of $\sigma_b(n)$ from a common profile.

          \item[(iv)] \textup{(Exponent concentration.)}
          O25~\cite{Beau2026a29} reports inter-pair standard deviations of the fitted pair exponent
          $0.539, 0.259, 0.161, 0.139, 0.109$ for $q = 29, 61, 101, 151, 211$, and $0.133, 0.214, 0.194$ for
          $q = 307, 401, 601$ from $24$, $12$ and $12$ sampled pairs (O25 Sections~3.2 and~3.4, Table~3).
          These concern fitted slopes of pair products, not the amplitudes $\sigma_b(n)$ block by block.

          \item[(v)] \textup{(No implication.)}
          Exact conjugation parity, identical pair exponents for all pairs and zero inter-pair spread are compatible
          with the failure of~\textup{[U]} for every $\varepsilon(q) < 1/3$.
        \end{enumerate}
      \end{proposition}

      \begin{proof}
        Items (i) to (iv) are statements about the cited sources at the locations given.
        For (v), let $f(n) = C\,(n+1)^{-\delta/2}$ on $\Nq$ and set $\sigma_b(n) = a_b\,f(n)$ with
        $a_b \in \{1, 2\}$, $a_{\bar b} = a_b$ for every conjugately matched pair, and both values occurring in $\Bq$.
        Parity holds exactly.
        Every pair product, for matched or independently sampled pairs alike, is $a_b a_{b'} f(n)^2$, whose
        logarithm is affine in $\log(n+1)$ with slope $-\delta$ on every subwindow, so all fitted pair exponents
        coincide, whatever the fitting window, and their spread is zero.
        If~\eqref{eq:universality} held with rate $\varepsilon$, a block with $a_b = 2$ and a block with $a_b = 1$
        would give $2f(n) \leq (1+\varepsilon)\sigma_*(n)$ and $f(n) \geq (1-\varepsilon)\sigma_*(n)$, hence
        $2(1-\varepsilon) \leq 1 + \varepsilon$, that is $\varepsilon \geq 1/3$.
        The data are abstract; the statement is one of non-implication, and it does not require them to be realised
        as rank-increment ratios.
      \end{proof}

      \begin{remark}[Scope of \textup{[U]}]\label{rem:U-scope}
        Hypothesis~[U] is open.
        U1~\cite{Beau2026u1} proves that no modulus of continuity in the reduced character can be inferred from the
        multiplication generators of the Heisenberg--Schr\"odinger representations, which are equidistant across
        distinct central characters,
        and locates a possible proof in the rank increment of the fingerprint filtration.
        Because $\sigma_b(n)$ is a ratio of integers with denominator $|S_n|$, a relative bound with small
        $\varepsilon$ on a fixed window is a combinatorial statement about Gram--Schmidt collisions.
        In O17's scalar model the raw count is exactly the same for every block, conjugate or not; uniformity of
        this kind can therefore hold for reasons that distinguish no spin sector, which is why [U] alone carries no
        geometric information and [R] is needed.
      \end{remark}

    \subsection{Conditions for exact conjugation parity}

      Write $\omega = e^{2\pi i/q}$ and let $\rho_c$ act on $\mathbb{C}^q$ by
      $[\rho_c(g_1, g_2, t)\psi](x) = \omega^{c(g_2 x+t)}\psi(x - g_1)$, the convention of O17.
      Then $\rho_{q-c}(g) = \overline{\rho_c(g)}$ entrywise in the standard basis, for every
      $g \in \Heis_3(\mathbb{Z}/q\mathbb{Z})$.

      \begin{lemma}[Conjugation transfer]\label{lem:conjugation}
        Let blocks $b$ and $\bar b$ have fingerprint families $\{v_b(g)\}$ and $\{v_{\bar b}(g)\}$ indexed by the same
        shells $S_0, S_1, \dots$, and assume
        \begin{equation}\label{eq:biparity}
          v_{\bar b}(g) = \lambda_g\,\overline{v_b(g)}, \qquad \lambda_g \in \mathbb{C}^\times,
        \end{equation}
        for every explored $g$ (conjugate matching), with the same normalisation $|S_n|$ for both blocks.
        Then $\delta r_n(\bar b) = \delta r_n(b)$ and $\sigma_{\bar b}(n) = \sigma_b(n)$ for every $n$.
        In particular, if $v_b(g) = \rho_c(g)\,w$ for a seed vector $w$, conjugate matching holds for the block with
        character $q-c$ and seed $\bar w$, with $\lambda_g = 1$.
      \end{lemma}

      \begin{proof}
        Entrywise conjugation $v \mapsto \bar v$ is a bijection of $\mathbb{C}^q$ that maps complex subspaces to
        complex subspaces of the same dimension and preserves linear dependence over $\mathbb{C}$; nonzero scalars
        change no span.
        Hence the span built over shells $0, \dots, n-1$ for $\bar b$ is the conjugate of that for $b$, and a vector
        of $S_n$ is independent of the first span exactly when its conjugate is independent of the second.
        Both readings of the increment (vector by vector against the previous span, or as a dimension increase) are
        therefore equal for $b$ and $\bar b$, and the common normalisation gives equality of $\sigma$.
        For the last assertion, $\rho_{q-c}(g)\,\bar w = \overline{\rho_c(g)}\,\bar w = \overline{\rho_c(g)\,w}$.
      \end{proof}

      Lemma~\ref{lem:conjugation} is exact but conditional.
      This paper does not verify that the O12/O13 construction of the fingerprint of a block
      $(c_1, c_2, c_3)$ has the seeded form, nor which block is conjugately matched to it.
      The O25 campaign samples the remaining entries of the two blocks of a pair independently, so its paired blocks
      are not selected to be conjugately matched and the lemma gives no pointwise equality for them.
      If a sampling rule maps each sampled block bijectively to a conjugately matched one and preserves the sampling
      weights, the lemma transfers to equality of the sampled statistics.
      Equality of observables is a spectral equivalence of blocks; it does not make a conjugate pair a fibre of the
      physical projection, which remains open (O18 Remark~3.2).

  \section{Stability of Ratios of Positive Averages}\label{sec:su2}

    \subsection{The abstract lemma}

      \begin{lemma}[Ratio stability]\label{lem:su2-isotropy}
        Let $I$ and $\Nq$ be finite sets, let $s \colon \Nq \to (0, \infty)$, and let $x_{i,n}$ satisfy
        $|x_{i,n} - s(n)| \leq \varepsilon\,s(n)$ for all $(i, n)$, with $0 \leq \varepsilon < 1$.
        Let $\alpha$ and $\alpha'$ be probability weights on $I$ and $\nu$ a probability weight on $\Nq$, and set
        \[
          a = \sum_{i,n} \alpha(i)\,\nu(n)\,x_{i,n}, \qquad
          a' = \sum_{i,n} \alpha'(i)\,\nu(n)\,x_{i,n}.
        \]
        Then $a, a' > 0$ and
        \begin{equation}\label{eq:ratio-bound}
          \left|\frac{a}{a'} - 1\right| \;\leq\; \frac{2\,\varepsilon}{1 - \varepsilon}.
        \end{equation}
        The bound is sharp: when the supports of $\alpha$ and $\alpha'$ are disjoint, equality holds for suitable data
        satisfying the hypothesis.
        If $\varepsilon \leq \frac12$, the right-hand side is at most $4\varepsilon$.
      \end{lemma}

      \begin{proof}
        Put $\bar s = \sum_n \nu(n)\,s(n) > 0$.
        Averaging the hypothesis with the weights $\alpha(i)\nu(n)$ gives
        $(1-\varepsilon)\bar s \leq a \leq (1+\varepsilon)\bar s$, and likewise for $a'$.
        Hence
        \[
          \frac{1-\varepsilon}{1+\varepsilon} - 1 \;\leq\; \frac{a}{a'} - 1
          \;\leq\; \frac{1+\varepsilon}{1-\varepsilon} - 1
          = \frac{2\varepsilon}{1-\varepsilon},
        \]
        and
        $1 - \frac{1-\varepsilon}{1+\varepsilon} = \frac{2\varepsilon}{1+\varepsilon}
        \leq \frac{2\varepsilon}{1-\varepsilon}$.
        Equality holds for $x_{i,n} = (1+\varepsilon)s(n)$ on the support of $\alpha$ and
        $x_{i,n} = (1-\varepsilon)s(n)$ on the support of $\alpha'$, when the two supports are disjoint.
        For $\varepsilon \leq \frac12$, $1 - \varepsilon \geq \frac12$.
      \end{proof}

      The common weight $\nu$ is used in the first line of the proof.
      With different depth weights $\nu \neq \nu'$, exact uniformity ($\varepsilon = 0$) gives
      $a/a' = \sum_n \nu(n)s(n) / \sum_n \nu'(n)s(n)$, which need not equal one.

      \begin{remark}[Role of conjugation parity]\label{rem:biparity-role}
        Split $\Cpm = \mathcal{C}_+ \sqcup \mathcal{C}_-$ according to the two nonzero weights.
        Under the hypotheses of Lemma~\ref{lem:conjugation}, and if a conjugate matching maps $\mathcal{C}_+$
        bijectively onto $\mathcal{C}_-$ and preserves the weights $\alpha_H$, the averages over
        $\mathcal{C}_+$ and over
        $\mathcal{C}_-$ agree exactly.
        This compares the two nonzero weights with each other; it says nothing about the comparison between the
        weight $\pm 1$ sector and the weight $0$ sector, which is the content of $A_H = A_Z$.
        The equality of the two horizontal coefficients in the real spatial form comes from the compatibility [INV]
        of Q7, not from parity.
      \end{remark}

    \subsection{Conditional coefficient form}

      \begin{proposition}[Coefficient isotropy under the inputs]\label{prop:conditional-isotropy}
        Assume \textup{[U]} with rate $\varepsilon(q) < 1$, \textup{[W]}, \textup{[R]} and \textup{[T]}.
        Then $A_H(q)/A_Z(q)$ satisfies~\eqref{eq:main}.
        Under \textup{[S]}, $|A_H(q) - a_q| \leq 2\,a_q\,\varepsilon(q)/(1-\varepsilon(q))$.
      \end{proposition}

      \begin{proof}
        Apply Lemma~\ref{lem:su2-isotropy} with $I = \Cpm \cup \Czero$, $x_{b,n} = \sigma_b(n)$, $s = \sigma_*$,
        $\alpha$ equal to $\alpha_H$ extended by zero and $\alpha'$ equal to $\alpha_Z$ extended by zero, and
        $\nu = \nu_q$; the hypothesis of the lemma is~\eqref{eq:universality}.
        By~[R] the ratio $a/a'$ equals $A_H(q)/A_Z(q)$, the coefficients of the supplied target [T].
        Under~[S], multiply by $A_Z(q) = a_q > 0$.
      \end{proof}

  \section{Invariant Forms and Proof of Theorem~\ref{thm:main}}\label{sec:casimir}

    \subsection{What invariance fixes}

      \begin{proposition}[Invariant forms on $\Sym^2(V_\rho)$]\label{prop:casimir-unique}
        Let $h_0$ be an $\mathfrak{su}(2)$-invariant Hermitian inner product on $\Sym^2(V_\rho)$.
        Every $\mathfrak{su}(2)$-invariant Hermitian form $h$ on $\Sym^2(V_\rho)$ is $h = a\,h_0$ for a real $a$, and
        $h$ is positive definite exactly when $a > 0$.
        For generators normalised by $[J_1, J_2] = i J_3$ and cyclically, the Casimir
        $\Csu = J_1^2 + J_2^2 + J_3^2$ acts as $2 \cdot \mathrm{Id}$, so $h_0(v, \Csu v) = 2\,h_0(v, v)$.
        The factor $a$ is not determined by invariance, and the eigenvalue $2$ depends on the normalisation of the
        generators.
      \end{proposition}

      \begin{proof}
        Write $h(v, w) = h_0(v, Tw)$ with $T$ self-adjoint for $h_0$.
        Invariance of $h$ and $h_0$ makes $T$ commute with the action, and $\Sym^2(V_\rho)$ is irreducible, so
        $T = a\,\mathrm{Id}$ by Schur's lemma; self-adjointness makes $a$ real.
        The Casimir acts on the spin-$j$ module as $j(j+1)$, which is $2$ for $j = 1$.
        Replacing $J_k$ by $\kappa J_k$ multiplies the Casimir by $\kappa^2$ and leaves the set of invariant forms
        unchanged.
      \end{proof}

      Consequently the positive invariant forms form one ray.
      The Casimir form, $v \mapsto h_0(v, \Csu v)$, is one point on it, singled out only by the normalisation of
      $h_0$ and of the generators.
      A transported spatial form that is $\mathfrak{su}(2)$-invariant would have $A_H = A_Z$, but its common value is
      an independent datum; this is the role of [S].
      The coincidence of a compression entry with the number $2$ is not normalisation-invariant and carries no
      evidential weight (Q7 Remark~6.3).

    \subsection{Proof of the main theorem}

      \begin{proof}[Proof of Theorem~\ref{thm:main}]
        The bound~\eqref{eq:main} and the bound under [S] are Proposition~\ref{prop:conditional-isotropy}.
        If $\varepsilon(q) \to 0$, the right-hand side of~\eqref{eq:main} tends to zero.
        If moreover $a_q \to a_* > 0$, then
        $|A_H(q) - a_*| \leq |A_H(q) - a_q| + |a_q - a_*| \to 0$.
        No input fixes $a_*$, and by Proposition~\ref{prop:casimir-unique} invariance does not fix it either.
      \end{proof}

      \begin{remark}[Epistemic status]\label{rem:epistemic}
        Lemma~\ref{lem:su2-isotropy}, Lemma~\ref{lem:conjugation} and Proposition~\ref{prop:casimir-unique} are
        unconditional mathematical statements.
        Proposition~\ref{prop:conditional-isotropy} and Theorem~\ref{thm:main} are conditional on [U], [W], [R] and
        [T], and their absolute part on [S].
        Hypothesis [U] is open (Remark~\ref{rem:U-scope}); its rate $\varepsilon(q)$ is not estimated, and no
        polynomial rate follows from any cited source.
        Inputs [W] and [R] have no construction in the corpus.
        Input [T] is the missing extended target of Q7 Conjecture~4.7.
        Input [S] has no source independent of $A_H$.
      \end{remark}

  \section{The Available Numbers}\label{sec:numerical}

    Table~\ref{tab:o25} lists, for the character $c = 7$ at the four O25 checkpoints, the stored mode index $k$, the
    corresponding diagonal entry of the three-row compression of the discrete Weil Laplacian computed in Q7 and its
    gap $\Delta$ to the flat-mode entry; the gap for $c = 3$; and, separately, the O25 inter-pair spread of fitted
    pair exponents, which concerns all sampled pairs and is a different quantity.

    \begin{table}[h]
      \centering
      \caption{Q7 compression entries for $c = 7$ and gaps for $c = 7$ and $c = 3$, and the O25 spread of fitted
      pair exponents over all sampled pairs, by prime $q$.}
      \label{tab:o25}
      \begin{tabular}{cccccc}
        \toprule
        $q$ & $k$ ($c=7$) & $2 + 4\sin^2(\pi k/q)$ & $\Delta$ ($c=7$) & $\Delta$ ($c=3$) & spread of $\dpair$ \\
        \midrule
        $61$  & $2$ & $2.042289$ & $0.042289$ & $0.094729$ & $0.259$ \\
        $101$ & $3$ & $2.034730$ & $0.034730$ & $0.095974$ & $0.161$ \\
        $151$ & $4$ & $2.027639$ & $0.027639$ & $0.084242$ & $0.139$ \\
        $211$ & $5$ & $2.022128$ & $0.022128$ & $0.088020$ & $0.109$ \\
        \bottomrule
      \end{tabular}
    \end{table}

    The columns derived from Q7 are identities, not measurements.
    The stored rows are pure Fourier modes of indices $\{0, -k, +k\}$, and each diagonal entry of the compression
    is the Rayleigh quotient $4 - 2\cos(2\pi m/q)$: exactly $2$ at the flat mode and $2 + 4\sin^2(\pi k/q)$ at the
    selected index (Q7 equations~(9)--(10) and Tables~1--2~\cite{Beau2026q7}).
    These entries are not coefficients of a continuum spatial form, and $\Delta$ is not an anisotropy.
    Its vanishing as $q \to \infty$ is equivalent to $\mathrm{dist}(k(q,c)/q, \mathbb{Z}) \to 0$, a property of
    the mode-selection rule (Q7 Conjecture~6.7); for $c = 3$ the sequence $\Delta(q)$ is not monotone.
    A power law fitted to the $c = 7$ column measures the growth of the selected index and supports no spectral
    rate.

    The last column is the concentration statistic of Proposition~\ref{prop:structural-U}(iv).
    It is not an estimate of $\varepsilon(q)$, which bounds amplitudes block by block rather than the spread of
    fitted slopes.
    Q5b Theorem~3.2 is a Gromov--Hausdorff convergence statement and carries no rate~\cite{Beau2026q5b}, so no
    exponent for $\varepsilon(q)$ or for $|A_H - A_Z|$ is predicted.

  \section{Consequences}\label{sec:consequences}

    \subsection{A conditional isotropic co-metric}

      \begin{corollary}[Conditional isotropic co-metric]\label{cor:metric}
        Assume the hypotheses of Theorem~\ref{thm:main} including \textup{[S]}, with $\varepsilon(q) \to 0$ and
        $a_q \to a_* > 0$.
        Assume further that the supplied operator $L_{\mathrm{ext}}$ of \textup{[T]} has temporal coefficient
        $-A_\tau$ with $A_\tau > 0$ and no mixed temporal--spatial terms in the adapted frame.
        Then its co-metric in that frame tends to $\mathrm{diag}(-A_\tau, a_*, a_*, a_*)$, of signature
        $(-,+,+,+)$.
      \end{corollary}

      \begin{proof}
        The spatial block is $\mathrm{diag}(A_H(q), A_H(q), A_Z(q))$ in the adapted frame by [T] and [INV], with
        $A_Z(q) = a_q$ by [S]; Theorem~\ref{thm:main} gives $A_H(q) \to a_*$.
        The temporal block is the stated assumption.
      \end{proof}

      No particular value of $a_*$, including $2$, is singled out; a value enters only through [S].
      The corollary takes $A_\tau$ independent of $q$.
      The Lorentzian signature of Q5b Theorem~6.1 concerns Q5b's own operator, conditional on Q5b Theorem~5.2,
      hence on [H-L] and [H-lift], and on a hyperbolicity hypothesis; it is not a statement about $L_{\mathrm{ext}}$.

    \subsection{Status of the bridge question}

      Theorem~\ref{thm:main} presupposes the extended target of Q7 Conjecture~4.7 through [T]; it does not construct
      it, and it does not establish that conjecture in any case, isotropic or not.
      Even when a supplied target is isotropic in the adapted frame, its equality with a transported Casimir form
      requires a scale to be specified (Q7 Remark~6.9).
      Problem Q5b-O2, the full-rank extension of Q5b's co-metric, is open, and Q8 shows that no lower-order
      correction of Q5b's operator produces a rank-three principal
      symbol or a constant positive central coefficient.

    \subsection{Open inputs}

      A derivation along this route needs, in order of logical dependence:
      \begin{enumerate}
        \item an extended spatial operator or effective-form construction supplying [T], with [ID], [ACT], [INV] and
        [ADAPT];
        \item a sector assignment [W] and a response law [R] linking that construction to fingerprint averages with a
        common depth weighting;
        \item a proof of [U] on the sector blocks, located in the rank increment of the fingerprint filtration;
        \item an absolute scale [S] obtained without reference to $A_H$;
        \item for exact parity, a verification that the production fingerprint of a block has the seeded form of
        Lemma~\ref{lem:conjugation}, together with a conjugately matched sampling rule.
      \end{enumerate}

  \section{Conclusion}\label{sec:conclusion}

    \subsection{What this paper establishes}

      Ratios of averages of uniformly controlled positive data with a common depth weighting are stable with the sharp
      bound $2\varepsilon/(1-\varepsilon)$.
      Exact conjugation parity of the fingerprint observable holds for conjugately matched data with an identical
      normalisation.
      Invariant Hermitian forms on the spin-one module form a single ray, so the Casimir eigenvalue fixes no
      coefficient.
      Under the named inputs [U], [W], [R] and [T], the coefficient ratio $A_H/A_Z$ tends to one; under [S] as well,
      $A_H$ tends to the supplied scale.
      None of these inputs is established, and the inputs cited in support of [U] do not imply it.

    \subsection{Interpretive outlook}

      \emph{This subsection is interpretation, not a result.}
      The structure of the argument separates two questions that a single number can appear to settle together.
      Relative isotropy, the equality of the horizontal and central coefficients, is the kind of property that could
      be inherited from the admissible fibre: if the geometric coefficients respond linearly to fingerprint data,
      then a fibre that looks the same from every character would produce a space that looks the same in every
      direction.
      The absolute scale is of a different kind.
      Representation theory fixes the shape of an invariant form and leaves its size free, so a physical scale for
      emergent space has to be supplied by some further datum of the substrate, such as a saturation bound, rather
      than read off a Casimir eigenvalue.
      For a reader situating the programme, the present status is that neither the linear response nor the
      extended spatial operator it presupposes has yet been constructed; the paper states what such constructions
      would have to deliver for isotropy to follow.

  \section*{Acknowledgements}

    Portions of the analysis and editorial development benefited from iterative interactions with large language
    models used as analytical assistants.
    All mathematical claims and interpretations remain the author's responsibility.

  \clearpage
  {\raggedright
  \bibliography{cosmochrony-bibliography}
  }

\end{document}
